\documentclass[../../main.tex]{subfiles}
\begin{document}

{\bf [解]}

$0 \sim 9$ の整数を足して $9$ になる2つずつで組をつくると以下の5組が存在する．
\begin{align}
  A: (0，9), B: (1，8), C: (2，7),
  D: (3，6), E:(4，5) \label{eq:1}
\end{align}
である．

\paragraph{(1)}

S のうち4桁の物を以下のように構成できる．
まず\cref{eq:1}の中から相異なる$4$つの組を選ぶ．
$4$つの組から数字を一つづつ選び，選んだ$4$つの数字を千の桁に$0$がこないように注意しながら配置した物である．

% TODO: 具体的な1つの要素の構成の仕方
\begin{figure}[htb]
  \centering
  \begin{tikzpicture}[scale=1.1, >=stealth, every node/.style={font=\small}]

    % ============================================================
    % ステップ 1: 5つのペアから相異なる4組を選ぶ（5C4）
    % ============================================================
    \node[font=\footnotesize\bfseries, blue!85!black, anchor=west] at (-0.5, 3.8)
    {\textbf{Step 1: 5組から4組のペアを選択} (${}_5\mathrm{C}_4 = 5$ 通り)};

    % 5つのペアのボックス配置
    % 未選択ペア A
    \node[draw=gray!50, fill=gray!10, rounded corners=3pt, inner sep=4pt, align=center] (pairA) at (0, 2.7)
    {\scriptsize $A$\\[1pt]$\{0, 9\}$};
    % 選択された4つのペア B, C, D, E (ハイライト)
    \node[draw=blue!85!black, fill=blue!10, very thick, rounded corners=3pt, inner sep=4pt, align=center] (pairB) at (2.2, 2.7)
    {\scriptsize $B$\\[1pt]$\{\mathbf{1}, 8\}$};
    \node[draw=blue!85!black, fill=blue!10, very thick, rounded corners=3pt, inner sep=4pt, align=center] (pairC) at (4.4, 2.7)
    {\scriptsize $C$\\[1pt]$\{2, \mathbf{7}\}$};
    \node[draw=blue!85!black, fill=blue!10, very thick, rounded corners=3pt, inner sep=4pt, align=center] (pairD) at (6.6, 2.7)
    {\scriptsize $D$\\[1pt]$\{\mathbf{3}, 6\}$};
    \node[draw=blue!85!black, fill=blue!10, very thick, rounded corners=3pt, inner sep=4pt, align=center] (pairE) at (8.8, 2.7)
    {\scriptsize $E$\\[1pt]$\{4, \mathbf{5}\}$};

    \node[font=\tiny, gray!60, below=1pt] at (pairA.south) {(未選択)};
    \draw[decorate, decoration={brace, amplitude=4pt}, thick, blue!85!black]
    (1.4, 3.3) -- (9.6, 3.3) node[midway, above=4pt, font=\scriptsize, blue!90!black] {選ばれた4組のペア};

    % ============================================================
    % ステップ 2: 各ペアから1つずつ数字を選ぶ（2^4）
    % ============================================================
    \node[font=\footnotesize\bfseries, teal!85!black, anchor=west] at (-0.5, 1.8)
    {\textbf{Step 2: 各ペアから数字を1つずつ選択} ($2^4 = 16$ 通り)};

    % 選出された数字のノード
    \node[draw=teal!85!black, fill=teal!15, very thick, circle, minimum size=0.75cm, font=\normalsize\bfseries] (num1) at (2.2, 0.9) {$1$};
    \node[draw=teal!85!black, fill=teal!15, very thick, circle, minimum size=0.75cm, font=\normalsize\bfseries] (num7) at (4.4, 0.9) {$7$};
    \node[draw=teal!85!black, fill=teal!15, very thick, circle, minimum size=0.75cm, font=\normalsize\bfseries] (num3) at (6.6, 0.9) {$3$};
    \node[draw=teal!85!black, fill=teal!15, very thick, circle, minimum size=0.75cm, font=\normalsize\bfseries] (num5) at (8.8, 0.9) {$5$};

    % 選択の矢印
    \draw[->, thick, teal!85!black] (pairB) -- (num1);
    \draw[->, thick, teal!85!black] (pairC) -- (num7);
    \draw[->, thick, teal!85!black] (pairD) -- (num3);
    \draw[->, thick, teal!85!black] (pairE) -- (num5);

    % ============================================================
    % ステップ 3: 4桁の各位に順列配置する（4!、ただし千の位≠0）
    % ============================================================
    \node[font=\footnotesize\bfseries, purple!85!black, anchor=west] at (-0.5, -0.1)
    {\textbf{Step 3: 各桁に配置して4桁の数を構成} ($4! = 24$ 通り，千の位 $\ne 0$)};

    % 4桁の位のボックス
    \coordinate (boxThous) at (2.7, -1.0);
    \coordinate (boxHund)  at (4.1, -1.0);
    \coordinate (boxTens)  at (5.5, -1.0);
    \coordinate (boxOnes)  at (6.9, -1.0);

    % 枠の描画
    \foreach \pos/\name/\digit in {
      boxThous/{千の位}/{$7$},
      boxHund/{百の位}/{$3$},
      boxTens/{十の位}/{$1$},
      boxOnes/{一の位}/{$5$}
    } {
      \node[draw=purple!85!black, fill=purple!8, very thick, rounded corners=2pt, minimum width=1.1cm, minimum height=1.1cm, font=\Large\bfseries]
      (\name) at (\pos) {\digit};
      \node[below=2pt, font=\scriptsize, gray!75!black] at (\name.south) {\name};
    }

    % 配置の対応矢印（例として 7 -> 千, 3 -> 百, 1 -> 十, 5 -> 一）
    \draw[->, dashed, thick, purple!75!black] (num7.south) to[out=250, in=90] (千の位.north);
    \draw[->, dashed, thick, purple!75!black] (num3.south) to[out=260, in=90] (百の位.north);
    \draw[->, dashed, thick, purple!75!black] (num1.south) to[out=290, in=100] (十の位.north);
    \draw[->, dashed, thick, purple!75!black] (num5.south) to[out=280, in=90] (一の位.north);

    % 完成した数値
    \node[draw=purple!90!black, fill=white, rounded corners=3pt, inner sep=4pt, font=\normalsize\bfseries, purple!95!black]
    at (9.3, -1.0) {構成例: $\mathbf{7315}$};

  \end{tikzpicture}
  \caption{4組のペアから相異なる数字を選び4桁の数を構成する手順}
  \label{fig:four_digit_construction_model}
\end{figure}

千の位に$0$に来てもいい時の場合の数は，
\begin{align}
  {}_5C_4 \times 2^4 \times 4! &= 1920 \text{個} \label{eq:2}
\end{align}
である．このうち千の位が0であるのは，
\begin{align}
  {}_4C_3 \times 2^3 \times 3! &= 192 \text{個} \label{eq:3}
\end{align}
だから，求める数は $1920 - 192 = 1728 \text{個}$ である．

\paragraph{(2)}

(1)と同様に数えると，
\begin{itemize}
  \item 1ケタのもの $\cdots$ 9個 \\
  \item 2ケタのもの $\cdots$ 72個 \\
  \item 3ケタのもの $\cdots$ 432個
\end{itemize}
だから，Sの要素を小さいものから並べた数列 $\{a_n\}$ において，4桁の数は $a_{514} \sim a_{2241}$ ，つまり$S_{2000}$ は 4桁の数である．

$S$の4桁の数のうち4ケタ目が $k$ ($1 \le k \le 9$) であるものは
\cref{eq:3}と同様に数えて 192個 あるから，
\begin{align}
  513 + (k-1) \cdot 192 < 2000 < 513 + k \cdot 192
\end{align}
をみたす $k \in \mathbb{N}$ は，$k=8$ であることとあわせて，$S_{2000}$ の千の位は 8 である．さらに，$S_{2000}$ は 千の位 が 8 であるもののうち，小さい方から数えて，
\begin{align}
  2000 - 1857 &= 143 \text{ 番目}
\end{align}
である．

次に百の位について千の位が$8$の時は\cref{eq:1}のうち残り A， C， D， E から 3つをえらぶことより，3位が $l$ ($l = 0,2,3,4,5,6,7,9$) である場合の数は，
\begin{align}
  {}_3P_2 \times 2^2 &= 24 \text{個}
\end{align}
である．よって，
\begin{align}
  24(m-1) < 143 \le 24m
\end{align}
をみたす $m \in \mathbb{N}$ は $m = 6$ だから，6番目の数字である $=6$ が百の位である．

$S_{2000}$ は $8 6 \square \square$ の形でかける $S$ の要素のうち大きい方から
\begin{align}
  144 - 143 + 1 = 2 \text{ 番目}
\end{align}
で，順に書き下すと，$8697$，$8695$
だから求めるのは
\begin{align}
  S_{2000} = 8695
\end{align}
である．

\end{document}
